\documentclass[12pt,a4paper]{article}
\usepackage[utf8]{inputenc}
\usepackage{amsmath, amssymb}
\usepackage{geometry}
\usepackage{setspace}
\usepackage{graphicx}

\geometry{margin=2.5cm}
\setstretch{1.2}

\title{Taller de Matemáticas Básicas II}
\author{José Luciano Mejía Arias}
\date{}

\makeatletter
\renewcommand{\maketitle}{
  \begin{flushright}
    \includegraphics[width=3cm]{../../logo_unal.png}
  \end{flushright}
  \begin{center}
    {\LARGE \textbf{\@title}}\\[0.3cm]
    {\large \@author}
  \end{center}
}
\makeatother

\begin{document}
\maketitle

\section*{Instrucciones}
Resuelva cada uno de los siguientes ejercicios mostrando todo el proceso. Presente su trabajo de manera ordenada y clara.

\section{Racionalización}
Racionalice el denominador de las siguientes expresiones:

\begin{enumerate}
    \item $\frac{5}{\sqrt{3}}$
    \item $\frac{2}{\sqrt{5} + 1}$
    \item $\frac{3}{\sqrt{7} - \sqrt{2}}$
    \item $\frac{\sqrt{2} + 1}{\sqrt{2} - 1}$
    \item $\frac{4}{\sqrt[3]{2}}$
    \item $\frac{x + \sqrt{x}}{x - \sqrt{x}}$
    \item $\frac{6}{2\sqrt{3} + 3\sqrt{2}}$
\end{enumerate}

\section{Inecuaciones}
Resuelva las siguientes inecuaciones y represente la solución en la recta numérica:

\begin{enumerate}
    \item $3x - 5 < 10$
    \item $2(x - 3) \geq x + 4$
    \item $-4 \leq 2x + 6 < 10$
    \item $x^2 - 5x + 6 > 0$
    \item $x^2 - 4x + 3 \leq 0$
    \item $\frac{x + 2}{x - 3} > 0$
    \item $|2x - 1| < 5$
    \item $\frac{2x - 1}{x + 3} \geq 1$
\end{enumerate}

\section{Conjuntos}
Dados los conjuntos $A = \{1, 2, 3, 4, 5\}$, $B = \{3, 4, 5, 6, 7\}$, $C = \{2, 4, 6, 8\}$ y $U = \{1, 2, 3, 4, 5, 6, 7, 8, 9, 10\}$ como conjunto universal, determine:

\begin{enumerate}
    \item $A \cup B$
    \item $A \cap B$
    \item $A - B$
    \item $B - A$
    \item $(A \cup B) \cap C$
    \item $A^c$ (complemento de A)
    \item $(A \cap B)^c$
    \item $A \times B$ (producto cartesiano, solo los primeros 10 elementos)
\end{enumerate}

\section{Operaciones con Conjuntos - Demostración}
Demuestre las siguientes propiedades usando diagramas de Venn o álgebra de conjuntos:

\begin{enumerate}
    \item $A \cup (B \cap C) = (A \cup B) \cap (A \cup C)$ (Ley distributiva)
    \item $(A \cup B)^c = A^c \cap B^c$ (Ley de De Morgan)
    \item $A - B = A \cap B^c$
\end{enumerate}

\section{Problemas de Aplicación con Conjuntos}

\begin{enumerate}
    \item En una encuesta a 100 estudiantes sobre sus preferencias deportivas, se encontró que:
    \begin{itemize}
        \item 45 practican fútbol
        \item 35 practican baloncesto
        \item 30 practican natación
        \item 15 practican fútbol y baloncesto
        \item 12 practican fútbol y natación
        \item 10 practican baloncesto y natación
        \item 5 practican los tres deportes
    \end{itemize}
    Determine:
    \begin{enumerate}
        \item ¿Cuántos estudiantes practican solo fútbol?
        \item ¿Cuántos estudiantes practican exactamente dos deportes?
        \item ¿Cuántos estudiantes no practican ninguno de estos deportes?
    \end{enumerate}
    
    \item En un grupo de 50 personas:
    \begin{itemize}
        \item 32 hablan inglés
        \item 28 hablan francés
        \item 25 hablan alemán
        \item 18 hablan inglés y francés
        \item 15 hablan inglés y alemán
        \item 14 hablan francés y alemán
        \item 10 hablan los tres idiomas
    \end{itemize}
    ¿Cuántas personas no hablan ninguno de estos tres idiomas?
\end{enumerate}

\section{Intervalos y Notación}
Exprese los siguientes conjuntos usando notación de intervalos y represéntelos en la recta numérica:

\begin{enumerate}
    \item $\{x \in \mathbb{R} : -3 < x \leq 5\}$
    \item $\{x \in \mathbb{R} : x \geq 2\}$
    \item $\{x \in \mathbb{R} : x < -1 \text{ o } x > 4\}$
    \item $\{x \in \mathbb{R} : |x - 2| < 3\}$
\end{enumerate}

\end{document}
